% ==========================================
% [BLOCO] DEDICATÓRIA
% ==========================================
\begin{dedicatoria}
   \vspace*{\fill}
   \begin{flushright}
   \noindent
   % [AGENT: INSERIR TEXTO DA DEDICATÓRIA AQUI]
   Este trabalho é dedicado aos meus pais, amigos e a todos 
   que contribuíram direta ou indiretamente para a minha formação.
   \end{flushright}
\end{dedicatoria}

% ==========================================
% [BLOCO] AGRADECIMENTOS
% ==========================================
\begin{agradecimentos}
% [AGENT: INSERIR TEXTO DOS AGRADECIMENTOS AQUI]
Agradeço primeiramente a Deus.
Aos meus familiares pelo apoio incondicional.
Ao meu orientador pela paciência e direcionamento durante toda a pesquisa.
\end{agradecimentos}

% ==========================================
% [BLOCO] EPÍGRAFE
% ==========================================
\begin{epigrafe}
    \vspace*{\fill}
    \begin{flushright}
    % [AGENT: INSERIR TEXTO DA EPÍGRAFE AQUI]
    \textit{``A imaginação é mais importante que o conhecimento.\\
    O conhecimento é limitado. A imaginação circunda o mundo.''\\
    (Albert Einstein)}
    \end{flushright}
\end{epigrafe}

% ==========================================
% [BLOCO] RESUMO EM PORTUGUÊS
% ==========================================
\begin{resumo}
    \noindent
    % [AGENT: INSERIR TEXTO DO RESUMO AQUI]
    O resumo deve ressaltar o objetivo, o método, os resultados e as conclusões do documento. A ordem e a extensão destes itens dependem do tipo de resumo (informativo ou indicativo) e do tratamento que cada item recebe no documento original.
    
    \vspace{\onelineskip}
    \noindent
    \textbf{Palavras-chave}: % [AGENT: INSERIR PALAVRAS-CHAVE AQUI SEPARADAS POR PONTO]
    TCC. LaTeX. Sistemas de Informação. IFS.
\end{resumo}

% ==========================================
% [BLOCO] ABSTRACT EM INGLÊS
% ==========================================
\begin{resumo}[Abstract]
    \begin{otherlanguage*}{english}
        \noindent
        % [AGENT: INSERIR TEXTO DO ABSTRACT AQUI]
        The abstract should highlight the objective, method, results, and conclusions of the document.
        
        \vspace{\onelineskip}
        \noindent
        \textbf{Keywords}: % [AGENT: INSERIR KEYWORDS AQUI SEPARADAS POR PONTO]
        Undergraduate Thesis. LaTeX. Information Systems. IFS.
    \end{otherlanguage*}
\end{resumo}
